\begin{textsample}{NEG Dim 9 – global\_north\_2024\_12 – Score -6.00 – global\_north\_2024\_12/118379632.txt}  \label{ex:f9_neg_124}
New \textbf{MEE} \textbf{newsletter} : Jerusalem Dispatch Sign up to get the latest insights and analysis on Israel-Palestine, alongside Turkey \textbf{Unpacked} and other \textbf{MEE} \textbf{newsletters} Mayors from northern towns flooded the media with threats to protest the decision.
Meanwhile, displaced Israelis said they couldn't return to their homes because the government had failed to deliver on its promise : total victory.
Little clarity Of the 70,000-plus Israelis who fled Israel ' s northern towns, few have gone back, with many claiming the security situation hadn't improved since Israel launched its attack against Hamas in Gaza, prompting an escalation of its conflict with Hezbollah, after the 7 October attacks on southern Israel. \textbf{MEE} asked some Israeli residents and analysts what a total victory would look like.
For them, one key demand was to have an official buffer zone in south Lebanon under Israeli security control, along with a disarmed Hezbollah.
For more than a year, Israeli Prime Minister Benjamin Netanyahu had claimed that Hezbollah intended to invade the Galilee.
Despite it being a far-fetched claim, for ordinary Israelis, an invasion of the Galilee is a highly plausible possibility as long as Hezbollah is armed and there ' s no buffer zone in place.
So why did Netanyahu agree to the ceasefire?
His televised address on 27 November offered little clarity.
Netanyahu tried to frame the deal as a temporary ceasefire, with some seeing it as a 60-day truce rather than an all-out end to the war.
That vagueness might have been deliberate as most of the residents in Israel ' s north and most of the local mayors are Likud voters who support Netanyahu.
New elites and the old guard Since the announcement, the Israeli media has been forced to ask why the country can sign an agreement with Lebanon but not one concerning Gaza."The new elites are interested in Gaza,"according to Rapoport,"so it was easy for Netanyahu to give up on Lebanon and concentrate on Gaza."According to Rapoport, the war on Lebanon was a war of Israel ' s old elite and its deep-state security establishment : a war of the Mossad and the Israel army, who are behind the assassination of Hezbollah leaders.
Meanwhile, the new elites, Finance Minister Bezalel Smotrich and National Security Minister Itamar Ben Gvir, are more interested in resettling Gaza.
Netanyahu, it seems, is mainly trying to preserve his coalition given that the ultra-Orthodox are threatening to ditch the government According to Rapoport, it looks like the Netanyahu government is more beholden to the new elites than the old guard.
Nehemia Shtrasler, a longtime Haaretz columnist, has also since analysed what likely led Netanyahu to sign the ceasefire agreement in Lebanon.
According to Shtrasler, the reasons Netanyahu gave in last week ' s speech were not the real ones.
Netanyahu, it seems, is mainly trying to preserve his coalition given that the ultra-Orthodox are threatening to ditch the government if the conscription exemptions they demand are not legislated soon.
So, the Lebanon ceasefire may have provided Netanyahu with some wiggle room on the domestic political front.
Netanyahu ' s ploy Yossi Verter of Haaretz highlighted the disparities : while Netanyahu gave the green light to a ceasefire with Lebanon just weeks after launching a ground operation, he is not even considering winding down the war in Gaza after 14 months.
Ben Caspit similarly asked in Maariv why it was it possible to make an agreement with Hezbollah, which is far better equipped to continue the fighting than Hamas, and why an agreement in Gaza seems somewhat impossible.
After mounting political losses, Israeli society may soon face a reckoning Yair Golan, heading the new Democrats Party ( a Meretz-Labor merger ), along with the families of the hostages, made similar remarks about the"missing"Gaza ceasefire.
While coalition interests may help shape Netanyahu ' s decisions about Gaza, Shtrasler claims the premier is also concerned about his personal safety.
The French foreign minister ' s promise not to enforce the ICC ' s arrest warrant suggests that signing the Lebanon ceasefire is a ploy to prevent his own arrest and outwit international sanctions.
In Tel Aviv, protests by the families of the 7 October abductees are still happening each week, prompting questions as to why Netanyahu has not shown the zeal for a deal to return the hostages as he has for a ceasefire on the northern front.
The bottom line is that the government doesn't appear to be in a hurry to end the war in Gaza.
Instead, what will most likely define the next phase of the broader conflict is Netanyahu ' s desire to warn Syria ' s Bashar al-Assad against aiding Hezbollah, and combating the threat posed by Iran.
The views expressed in this article belong to the author and do not \textbf{necessarily} reflect the editorial policy of Middle East Eye.
Lubna Masarwa is a journalist and Middle East Eye ' s Palestine and Israel bureau chief, based in Jerusalem.
Middle East Eye delivers independent and \textbf{unrivalled} coverage and analysis of the Middle East, North Africa and beyond.
To learn more about \textbf{republishing} this content and the associated fees, please fill out this form.
More about \textbf{MEE} can be found here.

% matched lemmas: mee, necessarily, newsletter, republish, unpack, unrivalled
\end{textsample}
